% ══════════════════════════════════════════════════════════════
%                          RÉSUMÉ
% ══════════════════════════════════════════════════════════════

\chapter*{Abstract}
\addcontentsline{toc}{chapter}{Abstract}

\vspace{0.3cm}

The management of internships in higher education remains largely fragmented and manual, particularly in Algeria, where students, companies, and universities rely on disconnected channels and paper-based workflows. This thesis presents the design and implementation of \textbf{Stag.io}, a modern full-stack web platform that centralizes and digitizes the complete internship lifecycle. The platform provides role-based dashboards for five types of actors (student, company, university, administrator, super administrator) and introduces \textbf{SmartMatch}, a proprietary recommendation engine that ranks internship offers against a student's profile across four weighted dimensions: skills (50\%), academic department (25\%), location (15\%), and title relevance (10\%). The system integrates real-time notifications via WebSocket, an AI-powered chatbot (Google Gemini), automated PDF generation for CVs and agreements, and full bilingual support (English/French). Built with Next.js, React, Express.js, PostgreSQL, and Drizzle ORM, the platform was validated through a comprehensive testing strategy covering the SmartMatch engine and security modules. Stag.io bridges the gap between academic institutions and the professional world, offering an efficient, transparent, and intelligent internship management experience.

\vspace{0.8cm}

\textbf{Keywords:} internship management, web platform, intelligent matching, recommendation engine, full-stack development, Next.js, React, PostgreSQL, WebSocket, UML.

\vspace{1.5cm}

\begin{center}
\rule{8cm}{0.5pt}
\end{center}

\vspace{0.5cm}

\selectlanguage{english}
\chapter*{Résumé}
\addcontentsline{toc}{chapter}{Résumé}

\vspace{0.3cm}

La gestion des stages dans l'enseignement supérieur reste largement fragmentée et manuelle, en particulier en Algérie, où les étudiants, les entreprises et les universités s'appuient sur des canaux déconnectés et des procédures papier. Ce mémoire présente la conception et la réalisation de \textbf{Stag.io}, une plateforme web full-stack moderne qui centralise et numérise l'ensemble du cycle de vie des stages. La plateforme offre des tableaux de bord adaptés à cinq types d'acteurs (étudiant, entreprise, université, administrateur, super administrateur) et introduit \textbf{SmartMatch}, un moteur de recommandation propriétaire qui classe les offres de stage par rapport au profil de l'étudiant selon quatre dimensions pondérées~: compétences (50\,\%), département académique (25\,\%), localisation (15\,\%) et pertinence du titre (10\,\%). Le système intègre des notifications en temps réel via WebSocket, un chatbot alimenté par l'IA (Google Gemini), la génération automatisée de CV et de conventions en PDF, ainsi qu'un support bilingue complet (anglais/français). Développée avec Next.js, React, Express.js, PostgreSQL et Drizzle ORM, la plateforme a été validée par une stratégie de test couvrant le moteur SmartMatch et les modules de sécurité.

\vspace{0.8cm}

\textbf{Mots-clés~:} gestion de stages, plateforme web, matching intelligent, moteur de recommandation, développement full-stack, Next.js, React, PostgreSQL, WebSocket, UML.
